% shared by main.tex (arXiv) and aee.tex (the anonymous submission); acmart hides the author under `anonymous`
% written by paper/macros.py — do not edit
\newcommand{\nEpisodes}{221{,}493}
\newcommand{\nModels}{19}
\newcommand{\nHypotheses}{23}
\newcommand{\nHolds}{6}
\newcommand{\nHolm}{6}
\newcommand{\replNineMT}{79\%}
\newcommand{\replNineShell}{97\%}
\newcommand{\replNineCode}{96\%}
\newcommand{\costNine}{3.5}
\newcommand{\replFourMT}{66\%}
\newcommand{\replFourShell}{97\%}
\newcommand{\replFourCode}{89\%}
\newcommand{\costFour}{7.4}
\newcommand{\replLfmTwoMT}{59\%}
\newcommand{\replLfmTwoShell}{50\%}
\newcommand{\replLfmTwoCode}{88\%}
\newcommand{\costLfmTwo}{11.1}
\newcommand{\replVlThirtyMT}{68\%}
\newcommand{\replVlThirtyShell}{59\%}
\newcommand{\replVlThirtyCode}{93\%}
\newcommand{\costVlThirty}{9.9}
\newcommand{\replLfmOneMT}{3\%}
\newcommand{\replLfmOneShell}{0\%}
\newcommand{\replLfmOneCode}{12\%}
\newcommand{\costLfmOne}{48.4}
\newcommand{\gemmaOld}{15\%}
\newcommand{\gemmaNew}{88\%}
\newcommand{\gemmaDiff}{$+73.8$}
\newcommand{\gemmaCI}{$[+66.5, +81.1]$}
\newcommand{\gemmaN}{64}
\newcommand{\qwenBrokenOld}{14.9\%}
\newcommand{\qwenBrokenNew}{0.6\%}
\newcommand{\memoryHelps}{2}
\newcommand{\instrNew}{$-1.1$}
\newcommand{\instrGateway}{$+6.0$}
\newcommand{\instrGap}{$-7.2$}
\newcommand{\instrGapCI}{$[-10.1, -4.2]$}
\newcommand{\switchFalse}{2\%}
\newcommand{\switchLo}{$+1.4$}
\newcommand{\switchHi}{$+2.1$}
\newcommand{\pbMT}{58\%}
\newcommand{\bigAloneQ}{80\%}
\newcommand{\tinyAloneMT}{0\%}
\newcommand{\tinyPlanMT}{43\%}
\newcommand{\fourAloneMT}{54\%}
\newcommand{\fourPlanMT}{56\%}
\newcommand{\bestVsPB}{$-2.8$}
\newcommand{\bestVsPBCI}{$[-8.5, +2.5]$}
\newcommand{\plannerShareLo}{96\%}
\newcommand{\feedbackSmall}{$+40.2$}
\newcommand{\feedbackFour}{$+57.7$}
\newcommand{\examplesSmall}{$+21.3$}
\newcommand{\examplesFour}{$+1.4$}
\newcommand{\levelLo}{55}
\newcommand{\levelHi}{59\%}
\newcommand{\protoLfm}{91\%}
\newcommand{\protoFour}{16\%}
\newcommand{\protoNine}{5\%}
\newcommand{\tinyStackQ}{0\%}
\newcommand{\liftTiny}{$+20.4$}
\newcommand{\liftTinyCI}{$[+16.8, +24.3]$}
\newcommand{\hurtFour}{$-15.9$}
\newcommand{\hurtFourCI}{$[-20.6, -11.2]$}
\newcommand{\hurtLfm}{$-5.9$}
\newcommand{\planCostFour}{2.5}
\newcommand{\routerVsPlan}{$+17.8$}
\newcommand{\routerVsPlanCI}{$[+13.2, +22.4]$}
\newcommand{\routerCost}{59\%}
\newcommand{\cpFourFalse}{14\%}
\newcommand{\tinyStackVsAlone}{$-37.9$}
\newcommand{\tinyStackVsAloneCI}{$[-42.1, -33.6]$}
\newcommand{\routerVsSingle}{$+3.1$}
\newcommand{\routerVsSingleCI}{$[+1.1, +5.4]$}

% written by paper/tables.py — do not edit
\newcommand{\verdictTable}{\begin{tabular}{@{}l p{0.62\textwidth} l r@{}}
\toprule
ID & Claim & Verdict & $p$ \\
\midrule
\multicolumn{4}{l}{\textit{Stage 1}} \\
A1 & Trajectory signals at step $\le$ 3 predict a small model's failure better than the prompt & not shown & -- \\
A2 & At least 40\% of small-model failures are recoverable obstacles & not shown & -- \\
A3 & A lever-matched cascade beats escalation alone at equal cost & holds & 0.0005 \\
A4 & The lift from a skill is $\cap$-shaped in model size & not shown & -- \\
A5 & A weight edit (WSP) lowers loops and raises quality & not shown & -- \\
\multicolumn{4}{l}{\textit{Stage 2}} \\
HR1 & Some harness component helps small models and not large ones, or the reverse & not shown & -- \\
HR2 & A harness chosen per model beats the best single harness & not shown & -- \\
HR3 & Retrying with the model's own harness beats a plain retry & holds & 0.0005 \\
HR4 & Routing (model, harness) pairs dominates routing models alone & holds & 0.0003 \\
HR5 & A rule that assigns harnesses is non-inferior to a search & not shown & -- \\
\multicolumn{4}{l}{\textit{Phase L}} \\
L1 & Worked examples help at least 5 of 7 new small models & not shown & -- \\
L2 & Instructions help at most 2 of 7, and less than large models & holds & 0.0001 \\
L3 & The corrected rule router is non-inferior to a search & not shown & -- \\
\multicolumn{4}{l}{\textit{Stage 3}} \\
PD1 & The best executor $\le$ 4B is within 5\,pp of the 27B executor & not shown & -- \\
PD2 & Delegation beats every single model at some budget & not shown & -- \\
PD3 & An executor $\le$ 1.5B is already within 5\,pp & not shown & -- \\
PD4 & Swapping the planner matters more than swapping the executor & not shown & -- \\
PD5 & Worked calls help both executors & not shown & -- \\
\multicolumn{4}{l}{\textit{Stage 4}} \\
CP1 & A 2.6B planner holds its role and beats its model alone & not shown & -- \\
CP2 & A 9B planner lifts a 350M executor & holds & 0.0005 \\
CP3 & Planning after a failure is non-inferior to planning always, at $\le$ 60\% of the cost & holds & 0.0005 \\
CP4 & The router beats the best single model at some budget & not shown & -- \\
CP5 & The router beats the tiny stack at the tiny stack's cost & not shown & -- \\
\bottomrule
\end{tabular}}
\newcommand{\parserTable}{\begin{tabular}{@{}l l r r r r r r@{}}
\toprule
Model & Old & Tasks & Old \% & Final \% & $\Delta$ pp & Broken old \% & Broken final \% \\
\midrule
gemma-4-26b-a4b & v1 & 64 & 14.6 & 88.4 & $+73.8$ & 0.3 & 0.0 \\
gemma-4-26b-a4b & v2 & 168 & 81.6 & 82.0 & $+0.4$ & 9.0 & 0.0 \\
gemma-4-26b-a4b & v3 & 1037 & 74.8 & 78.5 & $+3.6$ & 0.4 & 0.0 \\
granite-4.0-micro & v3 & 1347 & 37.8 & 40.0 & $+2.1$ & 1.9 & 1.9 \\
mnn-pocket-0.3b & v3 & 271 & 0.2 & 0.4 & $+0.2$ & 8.0 & 8.5 \\
qwen3-0.6b & v3 & 401 & 24.6 & 19.8 & $-4.8$ & 14.9 & 0.8 \\
qwen3-vl-30b & v3 & 881 & 81.9 & 74.9 & $-7.1$ & 0.7 & 0.0 \\
qwen3-vl-8b & v3 & 187 & 63.3 & 62.4 & $-0.9$ & 2.7 & 0.1 \\
qwen3.5-0.8b-wsp & v3 & 183 & 15.5 & 20.8 & $+5.3$ & 29.9 & 11.5 \\
qwen3.5-0.8b & v3 & 270 & 20.3 & 20.0 & $-0.3$ & 35.1 & 9.4 \\
qwen3.5-2b & v3 & 367 & 44.4 & 39.9 & $-4.5$ & 6.5 & 0.9 \\
qwen3.5-4b & v3 & 109 & 84.9 & 81.9 & $-3.0$ & 3.8 & 0.1 \\
qwen3.5-9b & v3 & 279 & 89.2 & 89.5 & $+0.4$ & 1.6 & 0.1 \\
qwen3.8 & v3 & 247 & 70.8 & 75.4 & $+4.6$ & 14.9 & 0.6 \\
\bottomrule
\end{tabular}}

\usepackage{booktabs}
\newcommand{\pp}{\,pp}
\setlength{\emergencystretch}{2em}  % let a line with a long \texttt name stretch rather than overflow
\settopmatter{printacmref=false}
\setcopyright{none}

\title{When Does a Small Model Pay?}
\subtitle{A Pre-Registered Study of Escalation, Harness Routing and Delegation for Agentic Small Language Models}

\author{Maxim Dunaev}
\orcid{0009-0009-2988-9839}
\affiliation{\institution{IDM Lab}\country{}}

\begin{abstract}
A recent position paper argues that small language models (SLMs) should be the default in agentic systems, with large
models called only when needed, but it reports no experiment. We test the claim on four agentic tracks with
\nModels{} open models from 0.23B to 27B active parameters and \nEpisodes{} logged episodes, in four stages that were
each pre-registered before their first episode. Every decision rule was also run on label-shuffled data, and a rule
that passed more than 5\% of shuffled copies did not count. Of \nHypotheses{} pre-registered hypotheses, \nHolds{} hold, and
all \nHolm{} survive a Holm correction across the study. Small models are cheap substitutes: a 4B model solves
\replFourMT{} of the multi-turn tasks a 27B model solves at a \costFour{}$\times$ lower cost per task, and replaceability
falls off below about 1.5B active parameters. Fixes in the harness are model-specific: of two findings from twelve
models, the one about instructions replicated on seven new models and the one about worked examples did not.
Delegation turns tiny models into useful executors (a 350M model goes from \tinyAloneMT{} to \tinyPlanMT{} of multi-turn
tasks with a 27B planner), but the planner is \plannerShareLo{} or more of the compute and the scheme does not pay.
Prompted small planners do not hold the role: a 2.6B planner breaks the protocol in \protoLfm{} of its turns, and a
tiny planner--executor stack solves \tinyStackQ{} of 650 tasks. A cheap 9B plan lifts a weak executor by \liftTiny\pp{}
but costs a strong one \hurtFour\pp{}; a router that plans only after the model fails alone beats planning always by
\routerVsPlan\pp{} at \routerCost{} of the cost, while its gain over the best single model did not survive calibration.
\end{abstract}
